\section{研究背景}

%%%%%%%%%%%%%%%
\begin{frame}{分布式应用}
  \fignocaption{width = 0.45\textwidth}{figures/sina-weibo-world-map.pdf}
  \vspace{0.50cm}

  \begin{columns}
	\column{0.40\textwidth}
	新浪微博社交应用~\footnotemark:
	\begin{itemize}
	  \item 日均用户近一亿名
	  \item 日均消息近一亿条
	\end{itemize}

	\column{0.50\textwidth}
	底层数据服务系统特性需求~\term{$H^{3}L$}: 
	\begin{itemize}
	  \item 低延迟, 高可用性 (4个9~\footnotemark)
	  \item 高容错性, 高可扩展性
	\end{itemize}
  \end{columns}
  
  \footnotetext[1]{2015第三季度; 数据来自 \href{https://www.chinainternetwatch.com/15740/weibo-q3-2015/}{China Internet Watch}.}
  \footnotetext[2]{数据来自 \href{http://www.infoq.com/cn/articles/weibo-platform-availability-9999}{InfoQ}.}
\end{frame}
%%%%%%%%%%%%%%%
\begin{frame}{分布数据}
  \graphicspath{{tikz-in-beamer/}}
  \begin{figure}[h!]
    \centering
    \begin{adjustbox}{max totalsize = {0.50\textwidth}{1.00\textheight}, center}
	  % File: beamer overlay for distributed-data.tex
% Created: July 23, 2016

\begin{tikzpicture}[rc/.style = {circle, minimum size = 4cm, fill = red},
  br/.style = {rectangle, rounded corners, minimum size = 4cm, fill = blue},
  gs/.style = {shape = star, scale = 8, star points = 5, star point ratio = 1.65, fill = green},
  dash pattern = {on 20pt off 10pt},
  partition/.style = {draw, lightgray, line width = 8pt},
  replication/.style = {draw, dash phase = 0pt, line width = 10pt, #1},
  access/.style = {draw, > = stealth, <->, dash phase = 0pt, line width = 6pt, brown}]

  \uncover<1->{
  \begin{pgfonlayer}{background}
	\node[opacity = 0.80] (map) {\includegraphics[scale = 0.50]{figs/world-map-logo.jpg}};
  \end{pgfonlayer}

  % shared data
  \begin{scope}[yshift = -18cm]
	\node[rc] (circ) {};
	\node[br, left = 2cm of circ] (rect) {};
	\node[gs, right = 2cm of circ] (st) {};

	\begin{pgfonlayer}{background}
	  \node[rectangle, rounded corners, fill = black!30, fit = (circ) (rect) (st), inner sep = 5pt, 
		label = {[font = \fontsize{65}{70}\sffamily\bfseries] below : Application Data}] (data) {};
	\end{pgfonlayer}
  \end{scope}
  }

  \uncover<2->{
  % partition
  \node[rc] (circ1) at (20, 10) {};
  \node[br] (rect1) at (-25, 10) {};
  \node[gs] (st1) at (28, -12) {};

  \draw[partition] (circ) to (circ1);
  \draw[partition] (rect) to (rect1);
  \draw[partition] (st) to (st1);
  }

  \uncover<3->{
  % replication
  \node[rc] (circ2) at (-15, -10) {};
  \node[rc] (circ3) at (-10, 18) {};

  \node[br] (rect2) at (10, 12) {};

  \node[gs] (st2) at (0,0) {};
  \node[gs] (st3) at (20, 5) {};

  \path[replication = {red}] (circ1) to (circ2) 
  	(circ2) to (circ3)
	(circ3) to (circ1);

  \path[replication = {blue}] (rect1) to (rect2);

  \path[replication = {green}] (st1) to (st2) 
  	(st2) to (st3)
	(st3) to (st1);
  }

%  \uncover<5->{
%  % sina-weibo clients
%  \def\SCALE{0.15}
%  \node[] (sina1) at (-25, 0) {\includegraphics[scale = \SCALE]{figs/sina-weibo.png}};
%  \node[] (sina2) at (10, 5) {\includegraphics[scale = \SCALE]{figs/sina-weibo.png}};
%  \node[] (sina3) at (5, -10) {\includegraphics[scale = \SCALE]{figs/sina-weibo.png}};
%
%  % access
%  \path[every edge/.append style = {access}] (sina1) edge (rect1)
%  				(sina1) edge (circ2)
%				(sina2) edge (rect2)
%				(sina2) edge (st2)
%				(sina3) edge (circ1)
%				(sina3) edge (st2);
%  }
\end{tikzpicture}

    \end{adjustbox}
  \end{figure}

  \only<1-3>{应用数据:} \only<4>{\textcolor{red}{分布数据 \term{distributed data}:}}
  \begin{enumerate}
	\item<2-> 分区 \term{partition}: 水平扩展
	\item<3-> 副本 \term{replication}: 就近访问, 容灾备份
  \end{enumerate}
\end{frame}
%%%%%%%%%%%%%%%
\begin{comment}
\begin{frame}{分布共享数据服务}
  \fignocaption{width = 0.55\textwidth}{figures/sina-weibo-redis.pdf}

  \begin{center}
	分布共享数据服务 \term{中间件}

	\term{Distributed Shared Data Service}
	\vspace{0.30cm}

	屏蔽底层数据分布性\; 提供共享数据抽象\; 简化上层应用开发
  \end{center}
\end{frame}
%%%%%%%%%%%%%%%
\begin{frame}{分布共享数据服务典型应用 (I)}
  \fig{width = 0.55\textwidth}{figures/dsss.pdf}
  {分布式存储系统 (\textcolor{blue}{\scriptsize 开源 [左]} \& \textcolor{red}{\scriptsize 商用 [右]}).}

  % \begin{description}
  %   \item[低延迟:] 就近访问副本数据
  %   \item[高可用性, 高容错性:] 备份容灾 
  % \end{description}
\end{frame}
%%%%%%%%%%%%%%%
\begin{frame}{分布共享数据服务典型应用 (II)}
  \fig{width = 0.75\textwidth}{figures/file-share.pdf}
  {个人多设备文件共享 {(\textcolor{blue}{\scriptsize [基于云] C/S 结构 [左]} \& 
  \textcolor{red}{\scriptsize P2P 结构 [右]}).}}

  \begin{description}
    \item[功能需求:] 文件副本 \citeinbeamer{Strauss}{MIT Thesis}{10}
    \item[网络断连:] 备份容灾; 离线可用
  \end{description}
\end{frame}
\end{comment}
%%%%%%%%%%%%%%%
% \begin{frame}{分布副本数据的典型应用 (三)}
%   \fig{width = 0.80\textwidth}{figures/coordination.pdf}
%   {分布式协同应用(上)及服务(下).}
% 
%   分布式协同服务需求 \citeinbeamer{Burrows}{OSDI}{06} vs. ``副本技术'':
%   \begin{description}
%     \item[高性能:] 低延迟就近``读''副本 \citeinbeamer{Yahoo!}{USENIXATC}{10}
%     \item[高可靠性:] 避免单点协同故障
%   \end{description}
% \end{frame}
%%%%%%%%%%%%%%%
